Abstract

This study proposed an offline reinforcement learning framework for tactical air combat.
We collected extensive offline datasets relevant to air combat and then adopted a decision transformer to leverage its sequential modeling ability. 
In addition, we proposed a lightweight maneuver pool to store and retrieve information for generating effective tactical maneuvers.
We believe that our study makes a significant contribution to the literature because the proposed algorithms can alleviate real-world trial-and-error costs.
Further, we believe that this paper will be of interest to the readership of your journal because the proposed framework is application-oriented and interdisciplinary.
This manuscript has not been published or presented elsewhere in part or entirety and is not under consideration by another journal.
We have read and understood your journal's policies, and we believe that neither the manuscript nor the study violates any of these.
There are no conflicts of interest to declare.

Dear Tamim Asfour,

We are pleased to submit our manuscript entitled "Maneuver-Conditioned Decision Transformer for Tactical In-Flight Decision-Making" for consideration in IEEE Robotics and Automation Letters (RA-L). 
This study introduces an innovative offline reinforcement learning framework tailored for tactical air combat. 
Our approach involves the collection of extensive offline datasets relevant to air combat scenarios, coupled with the adoption of a Decision Transformer to leverage its sequential modeling capabilities. 
Furthermore, we introduce a lightweight maneuver pool designed to store and retrieve information essential for generating effective tactical maneuvers.

We believe this study makes a significant contribution to the field, particularly in reducing trial-and-error costs associated with tactical air combat training and strategy development. 
The algorithms we propose offer practical solutions and advancements in the application of reinforcement learning in complex scenarios. 
Additionally, the interdisciplinary nature and application-oriented focus of our framework make this paper highly relevant to the readership of IEEE Robotics and Automation Letters (RA-L).

We confirm that this manuscript has not been published, or presented elsewhere in part or entirety, nor is it under consideration by another journal. We have thoroughly read and adhered to the journal's submission policies and guidelines. To the best of our knowledge, this manuscript does not violate any of these policies. There are no conflicts of interest to declare about this submission.

Thank you for considering our work for publication. We look forward to the possibility of contributing to IEEE Robotics and Automation Letters (RA-L) and eagerly await your feedback.

Sincerely,

Hoseong Jung

Agency for Defense Development, Daejeon, Republic of Korea
34186 Yuseong P.O.Box 35, Daejeon, Republic of Korea
E-mail: ghtjdaleka@gmail.com

\begin{thebibliography}{1}
    \bibitem{atari}
    V. Minh, K. Kavukcuoglu, D. Silver, A. Graves, I. Antonoglou, D. Wierstra, and M. Riedmiller, ``Playing atari with deep reinforcement learning,'' \emph{arXiv preprint arXiv:1312.5602}, 2013.

    \bibitem{shaw1985fighter}
    R. L. Shaw, ``Fighter combat,'' \emph{Tactics and Maneuvering: Naval Institute Press: Annapolis, MD, USA}, 1985.

    \bibitem{burgin1988rule}
    G. H. Burgin and L. Sidor, ``Rule-based air combat simulation,'' \emph{Tech. Rep.}, 1988.

    \bibitem{eklund2005implementing}
    J. Eklund, J. Sprinkle, and S. Sastry, ``Implementing and testing a non-linear model predictive tracking controller for aerial pursuit/evasion games on a fixed wing aircraft,'' in \emph{Proc. Am. Control. Conf.}, vol. 3, 2005, pp. 1509--1514.

    \bibitem{mcgrew2010aircombat}
    J. S. McGrew, J. P. How, B. Williams, and N. Roy, ``Air-combat strategy using approximate dynamic programming,'' \emph{J. Guid. Control. Dyn.}, vol. 33, no. 5, pp. 1641--1654, 2010.

    \bibitem{ernest2015genetic}
    N. Ernest, K. Cohen, E. Kivelevitch, C. Schumacher, and D. Casbeer, ``Genetic fuzzy trees and their application towards autonomous training and control of a squadron of unmanned combat aerial vehicles,'' \emph{Unmanned Syst.}, vol. 3, no. 3, pp. 185--204, 2015.

    \bibitem{shin2018autonomous}
    H. Shin, J. Lee, H. Kim, and D. H. Shim, ``An autonomous aerial combat framework for two-on-two engagements based on basic fighter maneuvers,'' \emph{Aerosp. Sci. Technol.}, vol. 72, pp. 305--315, 2018.

    \bibitem{yang2023manual}
    K. Yang, S. Kim, Y. Lee, C. Jang, and Y. -D. Kim, ``Manual-based automated maneuvering decisions for air-to-air combat,'' \emph{J. Aerosp. Inf. Syst.}, pp. 1--9, 2023.

    \bibitem{pope2022hierarchical}
    A. P. Pope, J. S. Ide, D. Mićović, H. Diaz,, J. C. Twedt, K. Alcedo, T. T. Walker, D. Rosenbluth, L. Ritholtz, and D. Javorsek, ``Hierarchical reinforcement learning for air combat at {DARPA}'s {AlphaDogfight Trials},'' \emph{IEEE Trans. Artif. Intell.}, pp. 1--15, 2022.

    \bibitem{chai2023hierarchical}
    J. Chai, W. Chen, Y. Zhu, Z. -X. Yao, and D. Zhao, ``A hierarchical deep reinforcement learning framework for 6-{DOF} {UCAV} air-to-air combat,'' \emph{IEEE Trans. Syst. Man. Cybern. Syst.}, 2023.

    \bibitem{li2023autonomous}
    B. Li, J. Huang, S. Bai, Z. Gan, S. Liang, N. Evgeny, and S. Yao, ``Autonnomous air combat decision-making of {UAV} based on parallel self-play reinforcement learning,'' \emph{CAAI Trans. Intell. Technol.}, vol. 8, no. 1, pp. 64--81, 2023.

    \bibitem{seong2023temp}
    H. Seong and D. H. Shim, ``Temp{F}user: Learning tactical and agile flight maneuvers in aerial dogfights using a long short-term temporal fusion transformer,'' \emph{arXiv preprint arXiv:2308.03257}, 2023.

    \bibitem{bae2023deep}
    J. Bae, H. Jung, S. Kim, S. Kim, and Y. -D. Kim, ``Deep reinforcement learning-based air-to-air combat maneuver generation in a realistic environment,'' \emph{IEEE ACCESS}, vol. 11, no. 1, pp. 26427--26440, 2023.

    \bibitem{yang2020maneuver}
    Q. Yang, J. Zhang, G. Shi, J. Hu, and Y. Wu, ``Maneuver decision of {UAV} in short-range air combat based on deep reinforcement learning,'' \emph{IEEE ACCESS}, vol. 8, pp. 363--378, 2020.

    \bibitem{prudencio2023survey}
    R. F. Prudencio, M. R. Maximo, and E. L. Colombini, ``A survey on offline reinforcement learning: Taxonomy, review, and open problems,'' \emph{IEEE Trans. Neural. Netw. Learn. Syst.}, 2023.

    \bibitem{nair2020awac}
    A. Nair, A. Gupta, M. Dalal, and S. Levine, ``{AWAC}: Accelerating online reinforcement learning with offline datasets," \emph{arXiv preprint arXiv:2006.09359}, 2020.

    \bibitem{chen2021decision}
    L. Chen, K. Lu, A. Rajeswaran, K. Lee, A. Grover, M. Laskin, P. Abbeel, A. Srinivas, and I. Mordatch, ``Decision Transformer: Reinforcement learning via sequence learning,'' \emph{Adv. Neural. Inform. Process. Syst.}, vol. 34, pp. 15084--15097, 2021.

    \bibitem{lee2022multi}
    K. -H. Lee, O. Nachum, M. Yang, L. Lee, D. Freeman, W. Xu, S. Guadarama, I. Fischer, E. Jang, H. Michalewski and I. Mordatch, ``Multi-game Decision Transformers,'' \emph{Adv. Neural. Inform. Process. Syst.}, vol. 35, pp. 27921--27936, 2022.

    \bibitem{fu2020d4rl}
    J. Fu, A. Kumar, O. Nachum, G. Tucker, and S. Levine, ``{D4RL}: Datasets for deep data-driven reinforcement learning,'' \emph{arXiv preprint arXiv:2004.07219}, 2020.

    \bibitem{berndt2004jsbsim}
    J. Berndt, "{JSBSim}: An open source flight dynamics model in {C}++,'' in \emph{AIAA Model. Simul. Technol. Conf. Exhib.}, 2004, p. 4923.

    \bibitem{haarnoja2018soft}
    T. Haarnoja, A. Zhou, P. Abbeel, and S. Levine, ``Soft actor-critic: Off-policy maximum entropy deep reinforcement learning with a stochastic actor,'' in \emph{Int. Conf. Mach. Learn.}, 2018, pp. 1861--1870.

    \bibitem{monastirsky2022learning}
    M. Monastirsky, O. Azulay, and A. Sintov, ``Learning to throw with a handful of samples using decision transformers,'' \emph{IEEE Robot. Autom. Letters},vol. 8, no. 2, pp. 576--583, 2022.

    \bibitem{xu2022prompting}
    M. Xu, Y. Shen, S. Zhang, Y. Lu, D. Zhao, J. Tenenbaum, and C. Gan, ``Prompting decision transformer for few-shot policy generalization,'' in \emph{Int. Conf. Mach. Learn.}, 2022, pp. 24631--24645.

    \bibitem{xu2022hyper}
    M. Xu, Y. Lu, Y. Shen, S. Zhang, D. Zhao, and C. Gan, ``Hyper-decision transformer for efficient online policy adaptation,'' in \emph{Int. Conf. Learn. Represent.}, 2022.

    \bibitem{liu2022few}
    H. Liu, D. Tam, M. Muqeeth, J. Motha, T. Huang, M. Bansal, and C. Raffel, ``Few-shot parameter-efficient fine-tuning is better and cheaper than in-context learning,'' \emph{Adv. Neural. Inform. Process. Syst.}, vol. 35, pp. 1950--1965, 2022.

    \bibitem{zhou2023real}
    G. Zhou, L. Ke, S. Srinivasa, A. Gupta, A. Rajeswaran, and V. Kumar, ``Real world offline reinforcement learning with realistic data source,'' in \emph{IEEE Int. Conf. Robot. Autom.}, 2023. pp. 7176--7183.
    
\end{thebibliography}
